\documentclass[10pt,a4paper]{amsart}
\usepackage[margin=19mm]{geometry}
\usepackage{amsmath,amssymb,mathtools}
\usepackage[scr=rsfs,cal=euler]{mathalfa}
\usepackage{xfrac}
\usepackage[unicode,hidelinks,bookmarks=false]{hyperref}
\newcommand{\ie}{i.e.\ }
\newcommand{\cf}{cf.\ }
\newcommand{\resp}{respectively\ }
\newcommand{\Subsection}{\subsection}
\providecommand{\nobreakdash}{\nobreak\hbox{-}}
\setlength{\emergencystretch}{2em}
\multlinegap0pt
\usepackage{algorithm}\usepackage{algpseudocode}
\usepackage{xcolor}
\usepackage{float}
\newtheorem{thm}{Theorem}
\title{Regularized coordinate minimization for nonconvex composite optimization with application to quantized image compression}
\author{Daniela Lupu, George T. Samoila, Adina M. Florea, Ion Necoara}
\date{Source: arXiv:2609.06260v1 (CC BY 4.0)}
\begin{document}
\maketitle

\noindent\emph{Source and licence.} Regularized coordinate minimization for nonconvex composite optimization with application to quantized image compression. Version arXiv:2609.06260v1 (CC BY 4.0), distributed by arXiv under the Creative Commons Attribution 4.0 International licence, http://creativecommons.org/licenses/by/4.0/, as recorded in the arXiv licence metadata for this exact version and shown on the abstract page https://arxiv.org/abs/2609.06260v1. Full licence notice: \texttt{licenses/arxiv-2609.06260-CC-BY-4.0.txt}.\par\smallskip

\noindent\textit{Editorial scope.} Counted unit: the article's thm thm:kl-decrease (source0.tex lines 304--318). The article's complete original proof of it is delivered with it (source0.tex lines 320--367, one \texttt{proof} environment of 48 lines). No mathematics is written, repaired or re-derived: the statement and the proof are the article's own lines, in the article's own notation, and no retained line is substituted or re-flowed.  Retained uncounted as the source-local context the proof needs: lines 298--301.  Nothing was omitted from any retained span, so the package carries no omission record.  No retained line contains a citation, so the wrapper carries no bibliography and no bibliography entry.  No retained line contains a figure environment, an image-inclusion command or drawing code, so no image is delivered and the package declares no asset dependency.  Editorial formatting only, in addition: an AMS \texttt{amsart} preamble that restates the article's own declarations and macros used by the retained lines, namely the environments \texttt{thm} on separate counters, together with the standard packages those lines need.  The retained environments print in this document as Theorem 1 in order of appearance, whereas the article prints the same items with the numbering of its own full document.  The complete native source of the article as distributed in the arXiv e-print for this exact version is bundled unchanged as \texttt{sources/209511/source0.tex}.
	\begin{align}
    \label{kl}
		& 	F(x) - F_*  \leq \sigma_q \text{dist}(0, \partial  F(x))^q,  
	\end{align}

\begin{thm}
\label{thm:kl-decrease}
Let Assumption~1 hold, and assume that the sequence of iterates $({x_k})_{k\geq 0}$ generated by algorithm~RCCM is bounded. Additionally, we assume that $F$ satisfies the KL property \eqref{kl} and  that there exists $k_0\geq 0$ such that, for every $k\geq k_0$, the iterates $x_{k}$ belong to this KL neighborhood. Then, for every $k\geq k_0$, we have:
\begin{align}
\label{eq:kl-decrease-final}
F(x_{k+1}) \!- \!F^*
\!\leq\!
(F(x_k)
-F^*)\! -\!
\frac{M_{\min}}{2C^2}
\left(\!
\frac{F(x_{k+1})-F_*}{\sigma_q}
\!\right)^{2/q}  .
\end{align}
\end{thm}

\begin{proof}
From the sufficient-decrease estimate obtained by combining
Lemma~1 and Lemma~3, we have:
\begin{align}
\label{eq:descent-kl-proof}
F(x_{k+1})
\leq
F(x_k)-
\frac{M_{\min}}{2C^2}
\min_{s_{k+1}\in\partial F(x_{k+1})}
|s_{k+1}|^2.
\end{align}
Since we have
$\operatorname{dist}(0,\partial F(x_{k+1}))
:=
\min_{s_{k+1}\in\partial F(x_{k+1})}
|s_{k+1}|$,
inequality \eqref{eq:descent-kl-proof} becomes
\begin{align}
F(x_{k+1})
\leq
F(x_k)
-
\frac{M_{\min}}{2C^2}
\operatorname{dist}(0,\partial F(x_{k+1}))^2.
\label{eq:descent-distance}
\end{align}
Further, for $k\geq k_0$, the iterate $x_{k}$ lies in the neighborhood
where the KL inequality \eqref{kl} holds, i.e.: 
\begin{align*}
F(x_{k+1})-F_*
\leq
\sigma_q
\left(\operatorname{dist}(0,\partial F(x_{k+1}))^2\right)^{\frac{q}{2}}.
\end{align*}
Substituting the above relation into \eqref{eq:descent-distance}, we obtain
\begin{align*}
F(x_{k+1})
\leq
F(x_k)
-
\frac{M_{\min}}{2C^2}
\left(
\frac{F(x_{k+1})-F_*}{\sigma_q}
\right)^{2/q},
\end{align*}
which proves \eqref{eq:kl-decrease-final}.
\end{proof}

\end{document}
